\section{Arquitectura del modelo: Qwen3-Next y Linear Attention}

\subsection{Arquitectura de atención híbrida}

Qwen3-Next adopta una \textbf{arquitectura Hybrid 3:1}: tres capas de Linear Attention por cada capa de Full Attention, apiladas de forma alternada. Este diseño y la arquitectura de Yuè Zhī Ànmiàn (月之暗面, Kimi) «llegan a lo mismo por caminos distintos».

\begin{importantbox}{¿Por qué cambiar de arquitectura?}
Hay tres motivaciones centrales:
\begin{enumerate}
    \item \textbf{Reducir el costo de KV Cache}: Linear Attention no necesita mantener un KV Cache completo, lo que ahorra memoria de la GPU de forma considerable
    \item \textbf{Dar soporte a contextos ultralargos}: prepara el terreno para infinite long context
    \item \textbf{Eficiencia de inferencia}: mayor velocidad de inferencia en escenarios de secuencias largas
\end{enumerate}
Además, el mecanismo Attention Gate del equipo de Qwen obtuvo el NeurIPS Best Paper.
\end{importantbox}

\subsection{Resumen de la sección}

La innovación arquitectónica es la infraestructura que sostiene las capacidades de la próxima generación (contextos ultralargos, interacción continua de los Agent). La arquitectura Hybrid Linear Attention de Qwen3-Next mantiene la calidad del modelo y, al mismo tiempo, ofrece la eficiencia necesaria para la interacción de Agent a largo plazo.
